% Beamer slides: Fortran Types, Expressions, and Control Flow

\documentclass[10pt,aspectratio=169]{beamer}
\usepackage[utf8]{inputenc}
\usepackage{listings}
\usepackage{hyperref}
\usepackage{xcolor}
\usepackage{ragged2e}
\usepackage{graphicx}
\usepackage{tikz}
\usetikzlibrary{positioning,calc,decorations.pathreplacing}

\usetheme{Dresden}
%\usetheme{Bergen}

\title{Programming Techniques 2026--2027}
\author{Hannu Parviainen}
\institute{Universidad de la Laguna}
\date{\today}

\lstset{
  language=fortran,
  basicstyle=\ttfamily\scriptsize,
  columns=fullflexible,
  keepspaces=true,
  backgroundcolor=\color{white},
  keywordstyle=\color{blue},
  commentstyle=\color{gray},
  stringstyle=\color{red},
  showstringspaces=false,
  breaklines=true,
  frame=lines
}

\title{Programming Techniques}
\subtitle{Fortran Types, Expressions, and Control Structures}
\author{Hannu Parviainen}
\date{\today}

\begin{document}

% Title
\begin{frame}
  \titlepage
\end{frame}

\begin{frame}[fragile]{Example: Triangle Area (Module + Program)}
\begin{lstlisting}[language=Fortran]
MODULE Triangle_Operations
    IMPLICIT NONE
    CONTAINS
    FUNCTION Area(x,y,z)
        REAL :: Area ! function type
        REAL, INTENT(IN) :: x, y, z
        REAL :: theta, height
        theta = ACOS((x**2 + y**2 - z**2) / (2.0*x*y))
        height = x * SIN(theta)
        Area = 0.5 * y * height
    END FUNCTION Area
END MODULE Triangle_Operations

PROGRAM Triangle
    USE Triangle_Operations
    IMPLICIT NONE
    REAL :: a, b, c
    PRINT *, 'Welcome, please enter the lengths of the 3 sides.'
    READ *, a, b, c
    PRINT *, "Triangle's area:", Area(a,b,c)
END PROGRAM Triangle
\end{lstlisting}
\end{frame}


\begin{frame}{Coding Style (Recommendations)}
\begin{block}{Conventions}
\begin{itemize}
\item Always use \texttt{IMPLICIT NONE}.
\item Consider capitalising keywords/intrinsics; keep user identifiers consistent.
\item Use small, consistent indentation (2–4 spaces); indent blocks and units.
\item Include unit names in corresponding \texttt{END} statements.
\end{itemize}
\end{block}
\begin{block}{Note}
To fit code on slides, I have relaxed some conventions in teaching materials.
\end{block}
\end{frame}


\begin{frame}[fragile]{Source Form (Free Form)}
\begin{block}{Rules}
\begin{itemize}
\item Up to 132 characters per line.
\item \texttt{!} introduces a comment.
\item \texttt{\&} continues a logical line.
\item \texttt{;} separates multiple statements on one line.
\item Whitespace is not significant.
\end{itemize}
\end{block}
\begin{block}{Example}
\begin{lstlisting}[language=Fortran]
PRINT *, "This line is continued " &
// "on the next line"; END ! end of program
\end{lstlisting}
\end{block}
\end{frame}

\section{Types and Declarations}

\begin{frame}{Intrinsic Types}
  \begin{block}{Default Intrinsic Types}
  \begin{itemize}
    \item \textbf{Character}: \texttt{CHARACTER :: sex}\; \texttt{CHARACTER(LEN=12) :: name}
    \item \textbf{Logical}: \texttt{LOGICAL :: wed}
    \item \textbf{Numeric}: \texttt{REAL :: height}, \texttt{INTEGER :: age}, \texttt{COMPLEX :: val}
  \end{itemize}
  \end{block}
\end{frame}

\begin{frame}{Literal Constants}
  \begin{block}{Examples}
  \begin{itemize}
    \item Integers: \texttt{12345}
    \item Reals: \texttt{1.0}, \texttt{-6.6E-06}
    \item Logical: \texttt{.FALSE.}, \texttt{.TRUE.}
    \item Character: \texttt{"Mau'dib"}, \texttt{'Mau''dib'} (escaped quote)
  \end{itemize}
  \end{block}
  \begin{block}{Notes}
  \begin{itemize}
    \item Only two logical values.
    \item Reals contain a decimal point (and may use exponential form).
    \item Character literals may be delimited by single or double quotes.
  \end{itemize}
  \end{block}
\end{frame}

\begin{frame}{Implicit Typing (Turn It Off)}
  \begin{block}{Rule}
  Without declarations, variables default to: \texttt{INTEGER} if starting with I–N; otherwise \texttt{REAL}. This is error-prone (or, more plainly, stupid).
  \end{block}
  \begin{block}{Best Practice}
  Always put \texttt{IMPLICIT NONE} immediately after any \texttt{USE} statements.
  \end{block}
  \begin{block}{Pitfall (Fixed Form)}
  In fixed form, \texttt{DO 30 I = 1.1000} can be parsed as assigning to a variable \texttt{DO30I} instead of a loop. Avoid by using free form and explicit declarations.
  \end{block}
\end{frame}

\begin{frame}[fragile]{Declarations (Numeric and Logical)}
  \begin{block}{General Syntax}
  \texttt{<type> [, <attributes>] :: <vars> [= <value>]}
  \end{block}
  \begin{block}{Examples}
\begin{lstlisting}[language=Fortran]
REAL :: x
INTEGER :: i, j
LOGICAL, POINTER :: ptr
REAL, DIMENSION(10,10) :: y, z
INTEGER :: k = 4
\end{lstlisting}
  \end{block}
  \begin{block}{Arrays}
  \texttt{DIMENSION(10,10)} declares a 10$\times$10 array.
  \end{block}
\end{frame}

\begin{frame}[fragile]{Symbolic Constants (Parameters)}
  \begin{block}{Definition}
  Symbolic constants are declared with the \texttt{PARAMETER} attribute.
  \end{block}
  \begin{block}{Examples}
\begin{lstlisting}[language=Fortran]
REAL, PARAMETER :: pi = 3.14159
CHARACTER(LEN=*), PARAMETER :: son = 'bart', dad = "Homer"
\end{lstlisting}
  \end{block}
  \begin{block}{When to Use}
  \begin{itemize}
    \item A variable should take only one value.
    \item To avoid “magic numbers”.
    \item For maintainability: values may change later.
  \end{itemize}
  \end{block}
\end{frame}

\begin{frame}[fragile]{Initialisation}
  \begin{block}{Rules}
  Initial values may use literals or other \texttt{PARAMETER}s.
  \end{block}
  \begin{block}{Examples}
\begin{lstlisting}[language=Fortran]
REAL :: x, y = 1.0D5
INTEGER :: i = 5, j = 100
CHARACTER(LEN=5) :: light = 'Amber'
CHARACTER(LEN=9) :: gumboot = 'Wellie'  ! right-padded with blanks
LOGICAL :: on = .TRUE., off = .FALSE.
REAL, PARAMETER :: pi = 3.141592, radius = 3.5
REAL :: circum = 2 * pi * radius
\end{lstlisting}
  \end{block}
  \begin{block}{Notes}
  A subset of intrinsics are allowed in initialisation expressions (e.g., \texttt{KIND}, \texttt{LEN}, \texttt{HUGE}, \texttt{TINY}, \texttt{EPSILON}, etc.).
  \end{block}
\end{frame}

\section{Expressions and Operators}

\begin{frame}{Expressions}
  \begin{block}{Types}
  Expressions can be numeric (e.g., \texttt{x+1}), character (\texttt{a//b}), or logical (\texttt{p .AND. q}).
  \end{block}
  \begin{block}{Use}
  Valid wherever an expression of the appropriate type is expected.
  \end{block}
\end{frame}

\begin{frame}[fragile]{Assignment}
  \begin{block}{Rule}
  The left-hand side is a variable/object; the right-hand side is an expression of the same type.
  \end{block}
  \begin{block}{Examples}
\begin{lstlisting}[language=Fortran]
a = b
c = SIN(.7) * 12.7      ! SIN uses radians
name = initials // surname
bool = (a .EQ. b .OR. c .NE. d)
\end{lstlisting}
  \end{block}
\end{frame}

\begin{frame}{Intrinsic Numeric Operations}
  \begin{block}{Operators}
  \begin{itemize}
    \item \texttt{**} exponentiation (right-associative): \texttt{10**2}
    \item \texttt{*}, \texttt{/} multiplication, division: \texttt{10*7/4}
    \item \texttt{+}, \texttt{-} unary/binary plus/minus: \texttt{10+7-4}, \texttt{-3}
  \end{itemize}
  \end{block}
  \begin{block}{Notes}
  RHS of \texttt{**} must be scalar. Operators apply to literals, scalars, and arrays (elemental semantics where applicable).
  \end{block}
\end{frame}

\begin{frame}[fragile]{Relational Operators}
  \begin{block}{Syntax}
  \texttt{.GT. (>)}, \texttt{.GE. (>=)}, \texttt{.LE. (<=)}, \texttt{.LT. (<)}, \texttt{.NE. (/=)}, \texttt{.EQ. (==)}
  \end{block}
  \begin{block}{Example}
\begin{lstlisting}[language=Fortran]
bool = (i .GT. j)
IF (i .EQ. j) c = d
\end{lstlisting}
  \end{block}
  \begin{block}{Reals}
  Avoid exact equality with reals; compare within a tolerance:
\begin{lstlisting}[language=Fortran]
REAL :: tol
tol = 1.0e-4
IF (ABS(a - b) .LT. tol) same = .TRUE.
\end{lstlisting}
  \end{block}
\end{frame}

\begin{frame}{Logical Operators}
  \begin{block}{Operators}
  \begin{itemize}
    \item \texttt{.NOT.} — logical negation
    \item \texttt{.AND.} — conjunction
    \item \texttt{.OR.} — disjunction
    \item \texttt{.EQV.} — equivalence (same truth value)
    \item \texttt{.NEQV.} — non-equivalence (different truth value)
  \end{itemize}
  \end{block}
  \begin{block}{Truth Table Highlights}
  \begin{itemize}
    \item \texttt{.NOT. .TRUE.} is \texttt{.FALSE.}
    \item \texttt{.TRUE. .AND. .FALSE.} is \texttt{.FALSE.}
    \item \texttt{.TRUE. .OR. .FALSE.} is \texttt{.TRUE.}
  \end{itemize}
  \end{block}
\end{frame}

\section{Control Flow}

\begin{frame}{Control Flow}
  \begin{block}{Overview}
    Control constructs allow the normal sequential order of execution to be changed. Fortran 90 supports:
    \begin{itemize}
      \item Conditional execution: \texttt{IF}, \texttt{IF ... THEN ... ELSE}
      \item Loops: \texttt{DO}, \texttt{DO WHILE}
      \item Multi-way choice: \texttt{SELECT CASE}
    \end{itemize}
  \end{block}
\end{frame}

\begin{frame}[fragile]{IF Statement}
  \begin{block}{Syntax}
    \texttt{IF (logical-expression) exec-stmt}
  \end{block}
  \begin{block}{Example}
\begin{lstlisting}[language=Fortran]
IF (x .GT. y) maxi = x
IF (a*b+c <= 47) boolie = .TRUE.
IF (i .NE. 0 .AND. j .NE. 0) k = 1/(i*j)
\end{lstlisting}
  \end{block}
\end{frame}

\begin{frame}[fragile]{Block IF}
  \begin{block}{Example}
\begin{lstlisting}[language=Fortran]
IF (i .EQ. 0) THEN
  PRINT*, "i is Zero"
ELSE IF (i .GT. 0) THEN
  PRINT*, "i is greater than Zero"
ELSE
  PRINT*, "i must be less than Zero"
END IF
\end{lstlisting}
  \end{block}
  \begin{block}{Notes}
    Both ELSE and ELSE IF are optional. Indentation improves readability.
  \end{block}
\end{frame}

\begin{frame}[fragile]{Conditional Exit}
  \begin{block}{Example}
\begin{lstlisting}[language=Fortran]
i = 0
DO
  i = i + 1
  IF (i .GT. 100) EXIT
  PRINT*, "i is", i
END DO
PRINT*, "Loop finished. i now equals", i
\end{lstlisting}
  \end{block}
  \begin{block}{Explanation}
    EXIT jumps out of the current DO loop.
  \end{block}
\end{frame}

\begin{frame}[fragile]{Conditional Cycle}
  \begin{block}{Example}
\begin{lstlisting}[language=Fortran]
i = 0
DO
  i = i + 1
  IF (i >= 50 .AND. i <= 59) CYCLE
  IF (i > 100) EXIT
  PRINT*, "i is", i
END DO
PRINT*, "Loop finished. i now equals", i
\end{lstlisting}
  \end{block}
  \begin{block}{Explanation}
    CYCLE skips to the next iteration of the innermost loop.
  \end{block}
\end{frame}

\begin{frame}[fragile]{Named and Nested Loops}
  \begin{block}{Example}
\begin{lstlisting}[language=Fortran]
outa: DO
  inna: DO
    IF (a .GT. b) EXIT outa
    IF (a .EQ. b) CYCLE outa
    IF (c .GT. d) EXIT inna
    IF (c .EQ. a) CYCLE
  END DO inna
END DO outa
\end{lstlisting}
  \end{block}
  \begin{block}{Notes}
    EXIT or CYCLE can target specific named loops.
  \end{block}
\end{frame}

\begin{frame}[fragile]{DO WHILE Loop}
  \begin{block}{Example}
\begin{lstlisting}[language=Fortran]
DO WHILE (a .EQ. b)
  ...
END DO
\end{lstlisting}
  \end{block}
  \begin{block}{Equivalent}
\begin{lstlisting}[language=Fortran]
DO
  IF (a .NE. b) EXIT
  ...
END DO
\end{lstlisting}
  \end{block}
\end{frame}

\begin{frame}[fragile]{Indexed DO Loops}
  \begin{block}{Example}
\begin{lstlisting}[language=Fortran]
DO i = 1, 100, 1
  ...
END DO
\end{lstlisting}
  \end{block}
  \begin{block}{Syntax}
    \texttt{DO var = expr1, expr2, [expr3]} with expr3 as step (default 1).
  \end{block}
\end{frame}

\begin{frame}[fragile]{Loop Variations}
  \begin{block}{Examples}
\begin{lstlisting}[language=Fortran]
DO i = 1, 30, 2     ! 15 iterations
DO j = 30, 1, -2    ! countdown
DO k = 30, 1, 2     ! zero-trip (skipped)
DO l = 1, 30        ! stride default = 1
\end{lstlisting}
  \end{block}
\end{frame}

\begin{frame}[fragile]{SELECT CASE}
	   \begin{block}{Example}
	   	\begin{lstlisting}[language=Fortran]
SELECT CASE (i)
CASE (3,5,7)
   PRINT*, "i is prime"
 CASE (10:)
   PRINT*, "i is > 10"
 CASE DEFAULT
   PRINT*, "i is not prime and is < 10"
 END SELECT
 \end{lstlisting}
   \end{block}
   \begin{block}{Notes}
     Cleaner and often more efficient than nested IFs.
   \end{block}
 \end{frame}

\section{Wrap-up}

\begin{frame}{Takeaways}
  \begin{block}{Essentials}
  \begin{itemize}
    \item Use \texttt{IMPLICIT NONE}; declare all variables.
    \item Prefer symbolic constants for fixed values.
    \item Be explicit with source form, types, and operators.
    \item Control flow constructs: IF, DO, DO WHILE, SELECT CASE.
  \end{itemize}
  \end{block}
\end{frame}

\end{document}
